\documentclass[12pt,a4paper]{article}

%% ─── Packages ───────────────────────────────────────────────────────────────
\usepackage[T1]{fontenc}
\usepackage[utf8]{inputenc}
\usepackage{lmodern}
\usepackage[margin=2.5cm]{geometry}
\usepackage{amsmath,amssymb,amsthm}
\usepackage{bm}
\usepackage{booktabs}
\usepackage{array}
\usepackage{multirow}
\usepackage{graphicx}
\usepackage{xcolor}
\usepackage{hyperref}
\usepackage{natbib}
\usepackage{microtype}
\usepackage{setspace}
\usepackage{caption}
\usepackage{subcaption}
\usepackage{algorithm}
\usepackage{algpseudocode}
\usepackage{tikz}
\usepackage{pgfplots}
\pgfplotsset{compat=1.18}
\usetikzlibrary{arrows.meta,positioning,shapes.geometric,fit,backgrounds,calc}

\hypersetup{
  colorlinks=true,
  linkcolor=blue!60!black,
  citecolor=blue!60!black,
  urlcolor=blue!60!black
}

\setstretch{1.15}

%% ─── Theorem environments ────────────────────────────────────────────────────
\newtheorem{definition}{Definition}[section]
\newtheorem{remark}{Remark}[section]

%% ─── Custom commands ─────────────────────────────────────────────────────────
\newcommand{\xeq}{x_{\mathrm{eq}}}
\newcommand{\Meq}{\mathcal{M}_{\mathrm{eq}}}
\newcommand{\SE}{\mathrm{SE}(3)}
\newcommand{\SOthree}{\mathrm{SO}(3)}
\newcommand{\Rthree}{\mathbb{R}^3}
\newcommand{\pdata}{p_{\mathrm{data}}}
\newcommand{\puncond}{p_{\mathrm{uncond}}}
\newcommand{\Loss}{\mathcal{L}}
\newcommand{\vhat}{\hat{v}}
\newcommand{\norm}[1]{\left\|#1\right\|}

%% ─── Title ───────────────────────────────────────────────────────────────────
\title{\Large\bfseries
  BBFlow: Learning Conformational Ensembles of Proteins Based on Backbone Geometry\\[0.4em]
  \large A Comprehensive Technical Review and Documentation}

\author{
  Nicolas Wolf$^{1,2,3}$ \quad Leif Seute$^{2,1,3}$ \quad Vsevolod Viliuga$^{4,1}$\\
  Simon Wagner$^3$ \quad Jan Stühmer$^{2,5}$ \quad Frauke Gräter$^{1,2,3}$\\[0.6em]
  \small $^1$ Max Planck Institute for Polymer Research, Mainz, Germany\\
  \small $^2$ Heidelberg Institute for Theoretical Studies, Heidelberg, Germany\\
  \small $^3$ IWR, Heidelberg University, Heidelberg, Germany\\
  \small $^4$ SciLifeLab and DBB, Stockholm University, Stockholm, Sweden\\
  \small $^5$ IAR, Karlsruhe Institute of Technology, Karlsruhe, Germany
}

\date{arXiv:2503.05738v2 \quad NeurIPS 2025}

%% ═══════════════════════════════════════════════════════════════════════════════
\begin{document}
\maketitle
\thispagestyle{empty}

\begin{abstract}
\noindent
Proteins do not exist as rigid, static entities; rather, their biological function is intimately governed
by the ensemble of conformations they adopt under physiological conditions. Molecular Dynamics
(MD) simulation is the established method for sampling such conformational ensembles from the
Boltzmann distribution, yet the computational cost of exhaustively exploring conformational space
often renders large-scale application infeasible. Recent deep generative approaches have sought to
emulate MD by learning from pre-computed trajectory data, but the dominant paradigm relies on
fine-tuning large pre-trained protein folding models (AlphaFold 2, ESMFold) and on evolutionary
sequence information in the form of Multiple Sequence Alignments (MSAs), limiting applicability
to natural proteins and incurring prohibitive inference costs.

In this document we provide a comprehensive technical exposition of \textbf{BBFlow}
(Backbone-geometry-conditioned Flow matching), a generative model that generates
MD-like backbone conformational ensembles conditioned solely on the equilibrium backbone
structure of the query protein, without recourse to evolutionary data or pre-trained folding weights.
BBFlow formulates ensemble generation as conditional structure generation on the Riemannian manifold
$\SE^N$, introduces a novel geometry-aware conditional prior distribution based on geodesic interpolation,
and employs an SE(3)-equivariant graph neural network with Clifford Frame Attention (CFA) as its
backbone architecture.  The model achieves state-of-the-art accuracy on the ATLAS benchmark at
inference speeds more than 40-fold faster than the best existing baseline, generalises to de~novo
proteins for which no evolutionary information exists, and is the first ensemble generation model
demonstrated on multi-chain protein complexes.

This document first introduces all necessary prerequisite concepts---protein biochemistry, Boltzmann
statistical mechanics, Molecular Dynamics, Riemannian geometry, equivariant neural networks, flow
matching, and attention mechanisms---before surveying the landscape of prior ensemble generation
methods and then providing a thorough derivation of every component of BBFlow, its training
procedure, evaluation protocol, experimental results, and ablation analysis.
\end{abstract}

\newpage
\tableofcontents
\newpage

%% ═══════════════════════════════════════════════════════════════════════════════
\section{Introduction}
\label{sec:introduction}

The determination of protein structure has been one of the defining scientific challenges of
the past century.  Since Kendrew and Perutz first resolved the X-ray crystal structure of
myoglobin and haemoglobin in the late 1950s, structural biology has sought to link the
three-dimensional arrangement of atoms to molecular function.  The advent of AlphaFold~2
\citep{jumper2021} and its contemporaries \citep{baek2021} has largely resolved the
\emph{static} structure prediction problem for single-domain proteins: given an amino-acid
sequence, an accurate equilibrium structure can now be predicted in seconds.

However, protein function does not reside in a single static conformation.  Enzymes
undergo induced-fit and conformational selection to bind substrates \citep{benkovic2008};
ion channels gate through large-scale conformational changes; allosteric proteins
communicate regulatory signals via concerted motions spanning tens of nanometres
\citep{bernetti2024}; intrinsically disordered regions adopt transient secondary structure
upon binding partners.  In each case, it is the \emph{ensemble} of accessible
conformations---and the free-energy differences between them---that encodes biological
function.

To characterise this conformational landscape, Molecular Dynamics (MD) simulation
integrates Newton's equations of motion for all atoms, sampling from the Boltzmann
distribution at a fixed temperature and pressure.  Although MD is physically rigorous,
covering the relevant state space is computationally expensive.  Biologically important
transitions occur on timescales from microseconds to seconds, while atomistic MD
integration steps are on the order of femtoseconds ($10^{-15}$~s).  Even short-timescale
simulations (hundreds of nanoseconds), sufficient to probe local flexibility and
thermal fluctuations relevant to ligand binding and allosteric regulation, require on
the order of thousands of CPU-hours per protein.

Recent years have seen growing interest in \emph{machine learning-based MD emulators}---
generative models trained on MD trajectory data that can sample new conformations at
inference without running a physical simulation.  These approaches promise inference
times orders of magnitude smaller than MD while maintaining distributional accuracy.
The dominant strategy has been to repurpose the powerful but expensive machinery of
protein folding models: AlphaFlow \citep{jing2024alphaflow} fine-tunes AlphaFold~2
on ensemble data, ConfDiff \citep{wang2024confdiff} uses AlphaFold~2's pre-trained
sequence representation as a conditioning signal, and BioEmu \citep{lewis2025bioemu}
builds an AlphaFold-like architecture trained on large heterogeneous trajectory data.
While accurate, these approaches inherit the limitations of their folding-model dependencies:
they are slow (due to MSA computation or Evoformer processing), potentially biased toward
natural protein families for which evolutionary data is abundant, and unable to handle
de~novo designed proteins.

\medskip
\noindent\textbf{This work.}
We present BBFlow, a generative model that entirely decouples conformational ensemble
generation from sequence-level modelling.  BBFlow:
\begin{enumerate}
  \item \textbf{Conditions on backbone geometry, not sequences.}  The equilibrium backbone
    structure $\xeq$ is encoded as pairwise distances and equivariant local directions---both
    computable from atomic coordinates without any database search.
  \item \textbf{Introduces a conditional prior.}  Rather than starting flow-matching integration
    from unconditioned noise, BBFlow samples its initial state $x_0$ from a conditional prior
    obtained by geodesic interpolation between noise and $\xeq$.  This geometrically biases
    the trajectory toward the correct fold from the outset.
  \item \textbf{Eliminates evolutionary information.}  No MSA, no protein language model embeddings,
    no pre-trained folding weights.  The model is trained from scratch on the ATLAS
    dataset~\citep{vandermeersche2024atlas} in 3 GPU-days.
  \item \textbf{Generalises to de~novo and multi-chain proteins.}  By dropping residue-index
    features and conditioning purely on geometry, BBFlow transfers to protein sequences with
    no evolutionary history and to multi-chain complexes never seen during training.
  \item \textbf{Achieves 40$\times$ speedup} over AlphaFlow-T at comparable accuracy, generating
    one conformation of a 302-residue protein in 0.8 seconds on an NVIDIA A100 GPU.
\end{enumerate}

The remainder of this document is organised as follows.  Section~\ref{sec:background}
develops all prerequisite concepts in detail.  Section~\ref{sec:related} surveys the landscape
of prior ensemble generation methods.  Section~\ref{sec:method} provides a full technical
derivation of BBFlow.  Section~\ref{sec:experiments} describes the experimental setup and
evaluation protocol.  Section~\ref{sec:results} reports quantitative results.  Section~\ref{sec:ablation}
presents ablation analyses.  Section~\ref{sec:discussion} discusses limitations and broader
implications, and Section~\ref{sec:conclusion} concludes.

%% ═══════════════════════════════════════════════════════════════════════════════
\section{Background and Prerequisite Concepts}
\label{sec:background}

This section develops, from first principles, the concepts required to understand BBFlow.
Readers familiar with any sub-topic may safely skip the relevant sub-section.

%% ───────────────────────────────────────────────────────────────────────────
\subsection{Protein Structure: From Sequence to Backbone Frames}
\label{sec:proteins}

A protein is a polypeptide chain of $N$ amino-acid residues, each selected from a standard
alphabet of 20 genetically encoded types.  Adjacent residues are connected by peptide bonds,
forming a repeating backbone of nitrogen (N), alpha-carbon (C$_\alpha$), and carbonyl-carbon
(C) atoms.  Side-chains attached to C$_\alpha$ differentiate the residue types and govern
physico-chemical properties.

\paragraph{Backbone dihedral angles.}
The local geometry of the backbone is parametrised by three dihedral angles per residue:
$\phi$ (rotation around N--C$_\alpha$), $\psi$ (rotation around C$_\alpha$--C), and $\omega$
(rotation around C--N, nearly fixed at $180°$ due to partial double-bond character).
The Ramachandran plot of allowed $(\phi,\psi)$ pairs defines regions corresponding to
secondary structures: $\alpha$-helices (near $\phi=-60°$, $\psi=-45°$) and
$\beta$-strands (near $\phi=-120°$, $\psi=+120°$).

\paragraph{Residue frames as $\SE$ elements.}
In the frame-based representation adopted by AlphaFold~2 \citep{jumper2021} and
all subsequent geometric models, each residue $i$ is assigned a local coordinate frame:
\begin{equation}
  T_i = (R_i, z_i) \in \SE,
\end{equation}
where $R_i \in \SOthree$ is a rotation matrix and $z_i \in \Rthree$ is the position of
C$_\alpha$.  The rotation is constructed from the N, C$_\alpha$, C backbone atoms of
residue $i$ using the Gram-Schmidt procedure, yielding a canonical right-handed frame
whose $x$-axis points from N toward C, and whose $y$-axis lies in the N--C$_\alpha$--C
plane.  The complete backbone of a protein is then represented as:
\begin{equation}
  \label{eq:backbone}
  x = (T_1, T_2, \ldots, T_N) \in \SE^N,
\end{equation}
a point on the $6N$-dimensional Riemannian product manifold $\mathcal{M} \equiv \SE^N$.

\paragraph{Why frames?}
The frame representation is particularly natural for machine learning because:
(i) it is coordinate-free---the relative geometry between two residues can be read off
from $T_i^{-1} T_j$ without reference to any global coordinate system;
(ii) it enables physically meaningful interpolation via geodesics on $\SE$; and
(iii) it provides a natural home for equivariant message-passing.

%% ───────────────────────────────────────────────────────────────────────────
\subsection{Conformational Ensembles and the Boltzmann Distribution}
\label{sec:boltzmann}

At thermodynamic equilibrium, the probability of observing a protein in conformation $x$ is
given by the Boltzmann distribution:
\begin{equation}
  \label{eq:boltzmann}
  p(x) \propto \exp\!\left(-\frac{E(x)}{k_B T}\right),
\end{equation}
where $E(x)$ is the potential energy of conformation $x$, $k_B$ is Boltzmann's constant, and
$T$ is the absolute temperature.  The partition function $Z = \int \exp(-E(x)/k_BT)\,dx$
normalises the distribution.

\begin{definition}[Conformational Ensemble]
  The \emph{conformational ensemble} of a protein at temperature $T$ is the set of
  conformations weighted by $p(x)$ in Equation~\eqref{eq:boltzmann}.  Thermodynamic
  quantities of interest---free energies, entropies, average properties---are expectations
  under this distribution.
\end{definition}

The free energy landscape $F(x) = -k_B T \ln p(x)$ governs kinetics as well as
thermodynamics: local minima are metastable states, and the barriers between them determine
transition rates (Arrhenius kinetics).  For proteins, low-energy basins correspond to
folded structures, partially unfolded intermediate states, and ligand-bound or unbound conformations.

\begin{remark}
  The Boltzmann distribution is defined over all atomic coordinates, including solvent.  In
  practice, explicit-solvent MD marginalises over solvent degrees of freedom, while
  implicit-solvent models approximate this marginalisation analytically.  BBFlow operates on
  backbone C$_\alpha$ coordinates after MD has already performed this marginalisation.
\end{remark}

%% ───────────────────────────────────────────────────────────────────────────
\subsection{Molecular Dynamics Simulation}
\label{sec:md}

MD numerically integrates Newton's equations of motion:
\begin{equation}
  \label{eq:newton}
  m_k \ddot{\mathbf{r}}_k = -\nabla_{\mathbf{r}_k} E(\mathbf{r}_1, \ldots, \mathbf{r}_n),
\end{equation}
for all $n$ atoms with masses $m_k$ and positions $\mathbf{r}_k$.  The potential energy $E$
is specified by a molecular force field---an empirically parameterised function combining
bonded terms (bonds, angles, dihedrals) and non-bonded terms (van der Waals, electrostatics).
Popular force fields for proteins include AMBER, CHARMM, GROMOS, and OPLS; the ATLAS
dataset~\citep{vandermeersche2024atlas} used to train BBFlow was generated with GROMACS
\citep{abraham2015gromacs} using the CHARMM36m force field with TIP3P water \citep{jorgensen1983}.

\paragraph{Time scales and the sampling problem.}
At a $2$-fs integration timestep, simulating $100$~ns of dynamics requires $5 \times 10^7$
integration steps, each evaluating forces for all atom pairs ($\mathcal{O}(n^2)$ without cutoffs,
$\mathcal{O}(n \log n)$ with particle-mesh Ewald).  For a typical soluble domain of $\sim 300$
residues ($\sim 5{,}000$ heavy atoms, $\sim 30{,}000$ total atoms with explicit solvent), this
amounts to wall-clock times on the order of $100$~CPU-hours per $100$~ns trajectory.  Three
$100$-ns trajectories per protein (as in ATLAS) thus cost approximately $300$~CPU-hours per
protein.

Worse, the timescales of biologically relevant transitions are often far longer than $100$~ns.
The ergodic hypothesis requires that an MD simulation runs long enough to visit all accessible
states in proportion to their Boltzmann weight.  For proteins with slow conformational
transitions (millisecond timescales), this is computationally prohibitive.  Even for short-timescale
dynamics, MD remains expensive when applied at scale (e.g.\ screening thousands of candidate
protein designs for desired flexibility).

\paragraph{What does short-timescale MD capture?}
Despite their limitations, $100$--$300$~ns simulations provide valuable information:
\begin{itemize}
  \item \textbf{Local flexibility}: side-chain rotamer distributions, loop mobility, B-factor
    analogues.
  \item \textbf{Secondary structure fluctuations}: fraying of helix termini, strand separation.
  \item \textbf{Binding-site plasticity}: relevant for drug design and protein--protein interaction
    prediction.
  \item \textbf{Allosteric communication}: correlated motions between distal sites.
\end{itemize}
These are precisely the properties that BBFlow aims to reproduce.

%% ───────────────────────────────────────────────────────────────────────────
\subsection{Lie Groups: \texorpdfstring{$\SOthree$}{SO(3)} and \texorpdfstring{$\SE$}{SE(3)}}
\label{sec:lie}

The special orthogonal group $\SOthree$ is the Lie group of $3 \times 3$ rotation matrices:
\begin{equation}
  \SOthree = \{R \in \mathbb{R}^{3\times 3} : R^T R = I,\; \det(R) = +1\}.
\end{equation}
The associated Lie algebra $\mathfrak{so}(3)$ consists of $3\times 3$ skew-symmetric matrices,
which can be identified with $\Rthree$ via the hat map $\hat{}: \Rthree \to \mathfrak{so}(3)$.

The special Euclidean group $\SE$ is the semidirect product of $\SOthree$ with $\Rthree$:
\begin{equation}
  \SE = \SOthree \ltimes \Rthree = \{(R, z) : R \in \SOthree,\; z \in \Rthree\},
\end{equation}
with group multiplication $(R_1, z_1) \cdot (R_2, z_2) = (R_1 R_2,\; z_1 + R_1 z_2)$.
$\SE$ represents rigid-body transformations of 3D space: rotations followed by translations.

\paragraph{Riemannian structure on $\SE^N$.}
A Riemannian metric on $\SE^N$ is required to define geodesics (shortest paths on the manifold)
and tangent vectors (velocities along those paths).  Following \citet{yim2023frameflow}, the
metric on $\SE^N$ is chosen as the product metric: the rotation component uses the bi-invariant
metric on $\SOthree$ (trace inner product on $\mathfrak{so}(3)$), and the translation component
uses the standard Euclidean metric.  The key property of this product metric is that geodesics
on $\SE^N$ decompose into independent rotation and translation geodesics for each residue.

\paragraph{Geodesics and the logarithm map.}
The geodesic from rotation $R_0$ to $R_1$ at time $t \in [0,1]$ is
$R_t = R_0 \exp(t \log(R_0^T R_1))$, where $\log: \SOthree \to \mathfrak{so}(3)$ is the
matrix logarithm.  The tangent vector at $R_t$ pointing toward $R_1$ is:
\begin{equation}
  \log_{R_t}(R_1) = R_t \log(R_t^T R_1),
\end{equation}
often written $\mathrm{Log}_{R_t}(R_1)$ in the context of Riemannian manifold flow matching.
The translation geodesic is simply the straight line $z_t = (1-t)z_0 + t z_1$, with constant
tangent vector $z_1 - z_0$.

%% ───────────────────────────────────────────────────────────────────────────
\subsection{Equivariance and Invariance in Neural Networks}
\label{sec:equivariance}

\begin{definition}[Equivariance and Invariance]
  Let $G$ be a group acting on input space $\mathcal{X}$ via $\rho_{\mathrm{in}}$ and on output
  space $\mathcal{Y}$ via $\rho_{\mathrm{out}}$.  A function $f: \mathcal{X} \to \mathcal{Y}$ is
  $G$-\emph{equivariant} if
  \begin{equation}
    f(\rho_{\mathrm{in}}(g) \cdot x) = \rho_{\mathrm{out}}(g) \cdot f(x) \quad \forall g \in G,\; x \in \mathcal{X}.
  \end{equation}
  If $\rho_{\mathrm{out}}(g) = \mathrm{id}$ for all $g$, then $f$ is $G$-\emph{invariant}.
\end{definition}

For molecular modelling, $G = \SE$ (rigid-body transformations of 3D space).  Physical
observables such as energy and forces must be invariant (or covariant) under global rotations
and translations---there is no privileged coordinate frame in space.  A neural network
that processes 3D atomic positions should therefore satisfy:
\begin{equation}
  f(R x + t) = f(x) \quad \text{(for scalar outputs),}
  \quad \text{or} \quad
  f(R x + t) = R f(x) \quad \text{(for vector outputs).}
\end{equation}

Achieving SE(3)-equivariance architecturally, rather than by data augmentation, was a
significant challenge in geometric deep learning.  Early approaches used handcrafted invariant
features (distances, angles).  More expressive methods include:
\begin{itemize}
  \item \textbf{Tensor Field Networks} \citep{thomas2018}: operate on spherical harmonic
    decompositions of 3D point clouds.
  \item \textbf{SE(3)-Transformers} \citep{fuchs2020}: extend the self-attention mechanism to
    SE(3)-equivariant features using spherical harmonics.
  \item \textbf{E(n) Equivariant GNNs (EGNN)} \citep{satorras2021}: use inter-atomic distances
    as invariant edge features in a message-passing network, achieving equivariance with
    a simplified architecture.
  \item \textbf{Frame averaging} \citep{puny2022}: guarantee invariance by averaging predictions
    over a canonical set of reference frames.
\end{itemize}
BBFlow uses the frame-based SE(3)-equivariant GNN architecture of GAFL \citep{wagner2024gafl},
in which geometric features are represented via projective geometric algebra (see
Section~\ref{sec:cfa}).

%% ───────────────────────────────────────────────────────────────────────────
\subsection{Transformer Self-Attention}
\label{sec:attention}

The transformer \citep{vaswani2017} architecture processes a sequence of $N$ token embeddings
$\{h_i\}_{i=1}^N \in \mathbb{R}^d$ via self-attention:
\begin{equation}
  \mathrm{Attention}(Q, K, V) = \mathrm{softmax}\!\left(\frac{QK^T}{\sqrt{d_k}}\right) V,
\end{equation}
where $Q = HW_Q$, $K = HW_K$, $V = HW_V$ are linear projections of the token embedding
matrix $H \in \mathbb{R}^{N \times d}$ onto query, key, and value spaces of dimension $d_k$.
The softmax normalises over the key dimension, producing attention weights that reflect
pairwise compatibility between queries and keys.

Intuitively, the query $q_i$ encodes ``what information residue $i$ is looking for'', the key
$k_j$ encodes ``what information residue $j$ offers'', and the value $v_j$ carries the
information that is aggregated if the attention weight $\alpha_{ij}$ is high.  Multi-head
attention runs $H_{\mathrm{head}}$ independent attention operations in parallel with different
learned projections, then concatenates and projects:
\begin{equation}
  \mathrm{MultiHead}(Q,K,V) = \mathrm{Concat}(\mathrm{head}_1, \ldots, \mathrm{head}_{H}) W_O.
\end{equation}

\paragraph{Invariant Point Attention (IPA).}
AlphaFold~2 \citep{jumper2021} introduced Invariant Point Attention (IPA) for protein structure
modules, extending scalar attention to geometric 3D queries and keys.  In IPA, each residue
learns 3D query points $\{q_i^{(p)}\}$ and key points $\{k_j^{(p)}\}$ that are defined in
local residue frames and then projected to global coordinates.  The attention logits include
an additional term based on the squared Euclidean distances between global query and key points:
\begin{equation}
  a_{ij} \propto \sum_p \norm{T_i(q_i^{(p)}) - T_j(k_j^{(p)})}^2,
\end{equation}
which is invariant to global rotations and translations because distances are invariant.  Output
points are rotated back into local frames, ensuring SE(3)-equivariance of the overall module.
IPA naturally processes both scalar sequence features (pair embeddings from Evoformer) and
the local 3D geometry of the backbone, making it the core architectural primitive for all
subsequent backbone generation models including BBFlow.

%% ───────────────────────────────────────────────────────────────────────────
\subsection{Clifford Frame Attention (CFA)}
\label{sec:cfa}

GAFL \citep{wagner2024gafl} extends IPA to Clifford Frame Attention (CFA), which represents
geometric features---points, lines, planes, and rigid frames---as elements of the projective
geometric algebra (PGA) $\mathbf{G}(\mathbb{R}^{3,0,1})$, also known as $\mathbb{P}(\mathbb{R}^*_{3,0,1})$.

\paragraph{Projective Geometric Algebra.}
PGA is a graded algebra in which the basic elements (blades) of grade $0, 1, 2, 3, 4$
represent scalars, planes, lines, points, and pseudoscalars, respectively.  The geometric
product $ab$ of two algebra elements is both associative and distributive and encodes
reflection, rotation, and translation as sandwich products: $p' = m p \tilde{m}$, where
$m$ is a motor (an element encoding a screw motion) and $\tilde{m}$ its reverse.

\paragraph{CFA construction.}
In CFA, residue frames $T_i$ are represented as motors $M_i \in \mathbf{G}$.  Edge features
between residues $i$ and $j$ are encoded as multivectors $\mathbf{e}_{ij} \in \mathbf{G}$
using the relative motor $M_i^{-1} M_j$ and the pairwise geometry.  Messages are constructed
using bilinear products in the algebra (outer product, inner product, geometric product),
enabling higher-order geometric interactions than scalar inner products alone.  CFA thus
generalises IPA by leveraging the full algebraic structure of rigid-body geometry, leading
to more expressive geometric representations at comparable parameter count.

%% ───────────────────────────────────────────────────────────────────────────
\subsection{Diffusion Models}
\label{sec:diffusion}

Diffusion models \citep{song2021} define a forward process that gradually adds Gaussian noise
to data over $T$ timesteps:
\begin{equation}
  q(x_t | x_0) = \mathcal{N}(x_t; \sqrt{\bar{\alpha}_t} x_0, (1-\bar{\alpha}_t) I),
\end{equation}
and learn a reverse denoising network $\epsilon_\theta(x_t, t)$ that predicts the added noise.
Sampling is performed by iteratively denoising from pure Gaussian noise $x_T \sim \mathcal{N}(0,I)$:
\begin{equation}
  x_{t-1} = \frac{1}{\sqrt{\alpha_t}}\left(x_t - \frac{1-\alpha_t}{\sqrt{1-\bar{\alpha}_t}} \epsilon_\theta(x_t,t)\right) + \sigma_t z,
  \quad z \sim \mathcal{N}(0,I).
\end{equation}
The stochastic nature of the reverse process (Langevin dynamics) requires many ($\sim\!1000$)
denoising steps for high-quality generation, though DDIM \citep{song2021} and related
deterministic samplers reduce this to $\sim\!50$--$100$ steps.

A critical limitation of diffusion models on curved spaces (such as $\SOthree$ or $\SE^N$) is
that the Gaussian noise prior is not well-defined on non-Euclidean manifolds.  Riemannian
diffusion \citep{yim2023framediff} addresses this by defining the forward process via the
heat kernel on the manifold, but the Gaussianity of the prior remains a theoretical requirement.

%% ───────────────────────────────────────────────────────────────────────────
\subsection{Flow Matching}
\label{sec:flow_matching}

Flow matching \citep{lipman2023,chen2024} offers a more flexible and often faster-converging
alternative to diffusion.  The core idea is to learn a time-dependent vector field
$v_\theta: \mathcal{M} \times [0,1] \to T\mathcal{M}$ that transports samples from a simple
prior $p_0$ to the data distribution $p_1$ via an ODE:
\begin{equation}
  \label{eq:flow_ode}
  \frac{dx_t}{dt} = v_\theta(x_t, t), \qquad x_0 \sim p_0.
\end{equation}
The \emph{flow} $\phi_t: \mathcal{M} \to \mathcal{M}$ is defined implicitly as the solution map
of this ODE: $x_t = \phi_t(x_0)$.

\paragraph{Training via conditional flow matching.}
Direct regression on the marginal vector field $v(x,t)$ is intractable because it requires
the marginal density $p_t(x)$.  Conditional flow matching sidesteps this by connecting
\emph{individual} data pairs $(x_0, x_1)$ via a simple conditional path:
\begin{equation}
  \psi(x_0, x_1, t) = (1-t) x_0 + t x_1 \quad \text{(Euclidean case)},
\end{equation}
with conditional vector field $u_t = x_1 - x_0$.  The flow matching loss is:
\begin{equation}
  \label{eq:fm_loss}
  \Loss_{\mathrm{FM}} = \mathbb{E}_{t, x_0, x_1}\!\left[\norm{v_\theta(x_t, t) - u_t}^2\right],
  \quad x_t = \psi(x_0, x_1, t).
\end{equation}
It can be shown \citep{lipman2023} that the conditional loss~\eqref{eq:fm_loss} and the
marginal loss have the same gradient with respect to $\theta$, so minimising~\eqref{eq:fm_loss}
yields the correct marginal vector field.  This is the key insight that makes flow matching
tractable.

\paragraph{Advantages over diffusion.}
\begin{itemize}
  \item \textbf{Straight paths}: the interpolation $\psi$ is a straight line, producing nearly
    straight flow trajectories that can be integrated with very few (e.g.\ 20) ODE steps.
  \item \textbf{Flexible priors}: unlike diffusion, flow matching does not require a Gaussian
    prior.  Any distribution $p_0$ can be used, enabling the conditional prior of BBFlow.
  \item \textbf{Riemannian generalisation}: on manifolds $\mathcal{M}$, straight lines are
    replaced by geodesics $\gamma(x_0, x_1, t)$ \citep{chen2024}, and the conditional
    vector field is the geodesic velocity.
\end{itemize}

\paragraph{Flow matching on $\SE^N$.}
For the protein backbone manifold $\mathcal{M} = \SE^N$, the geodesic decomposes per-residue
into rotation and translation parts.  The ground-truth conditional vector fields are
\citep{yim2023frameflow}:
\begin{align}
  \label{eq:vfield_so3}
  v_{\SOthree}(R_t, t | R_1) &= \frac{\log_{R_t}(R_1)}{1-t}, \\
  \label{eq:vfield_r3}
  v_{\Rthree}(z_t, t | z_1) &= \frac{z_1 - z_t}{1-t}.
\end{align}
The neural network parametrises $x_1 = \hat{x}_\theta(x_t, t, \xeq)$, from which the
predicted vector fields $\hat{v}_{\SOthree}$ and $\hat{v}_{\Rthree}$ are recovered using
Equations~\eqref{eq:vfield_so3} and~\eqref{eq:vfield_r3}.  This parametrisation is analogous
to the $x_0$-prediction in diffusion models and empirically leads to faster convergence.

%% ═══════════════════════════════════════════════════════════════════════════════
\section{Prior Work: Approaches to Conformational Ensemble Generation}
\label{sec:related}

We organise the landscape of prior ensemble generation methods by their architectural strategy
and the degree to which they depend on evolutionary information or pre-trained folding models.

%% ───────────────────────────────────────────────────────────────────────────
\subsection{Classical Simulation-Based Methods}
\label{sec:classical}

\paragraph{Molecular Dynamics.}
As discussed in Section~\ref{sec:md}, MD is the gold-standard method for generating
conformational ensembles.  It is physically rigorous, requires no training data, and can
in principle sample any region of conformational space.  Its primary limitation is
computational cost: $\mathcal{O}(10^4)$ to $\mathcal{O}(10^6)$ CPU-hours per protein for
trajectories long enough to capture biologically relevant transitions.

\paragraph{Normal Mode Analysis (NMA).}
NMA \citep{case1994} approximates the potential energy surface near the equilibrium structure
as a harmonic (quadratic) function, then analytically computes the normal modes---the
orthogonal directions of concerted motion with definite frequencies.  NMA is extremely fast
(minutes per protein) but is a linear approximation that cannot capture large-amplitude
or anharmonic motions.  It also requires a minimised equilibrium structure.  Benchmarks in
\citet{jing2024alphaflow} show that NMA systematically underestimates ensemble spread
compared to MD and all generative baselines.

\paragraph{MSA Subsampling / AlphaFold2 Sampling.}
\citet{wayment2024} showed that by running AlphaFold~2 with different MSA subsets (each
encoding a different evolutionary sub-population), one can coax the model into predicting
alternative conformations.  This approach is powerful for proteins with clearly separable
conformational substates encoded in evolutionary variation, but is limited to proteins with
rich MSA coverage and cannot sample thermally accessible states not reflected in sequence
variation.

%% ───────────────────────────────────────────────────────────────────────────
\subsection{Boltzmann Generators}
\label{sec:boltzmann_gen}

\citet{noe2019} introduced Boltzmann Generators (BG), which train normalising flows to map
a simple Gaussian latent space to the Boltzmann distribution $p(x) \propto e^{-E(x)/k_BT}$.
The training uses a combination of maximum likelihood on MD samples (density estimation)
and physics-based energy minimisation (the free-energy loss).

\paragraph{Strengths.}
BGs can sample thermodynamically equilibrated configurations, including rare states not
well-sampled by short MD, because they directly target the Boltzmann distribution.  They
require relatively few MD samples for training.

\paragraph{Limitations.}
Crucially, BGs must be trained \emph{per protein}: the normalising flow is fitted to the
specific energy landscape of a single molecule.  They are not \emph{transferable}---a model
trained on one protein cannot generate ensembles for a new protein without retraining.  This
severely limits their practical utility for large-scale screening.  More recent work on
equivariant BGs \citep{klein2023} has partially addressed the symmetry constraints but the
transferability issue remains.

%% ───────────────────────────────────────────────────────────────────────────
\subsection{AlphaFlow: AlphaFold-Based Ensemble Generation}
\label{sec:alphaflow}

\citet{jing2024alphaflow} proposed AlphaFlow, the first transferable protein ensemble generation
model, by reformulating ensemble generation as a \emph{structure denoising} task within the
AlphaFold~2 framework.  Three variants were introduced:

\paragraph{AlphaFlow (no templates).}
The pre-trained AlphaFold~2 model (Evoformer + structure module) is fine-tuned on MD trajectory
data from ATLAS.  At inference, the model predicts a protein conformation from its sequence and
MSA, starting from a noisy structure.  The flow matching process runs over the AlphaFold structure
module while the Evoformer processes the MSA once per sequence.

\paragraph{AlphaFlow-T (with templates).}
The equilibrium structure $\xeq$ is provided as a \emph{template} to AlphaFold's template
processing module, conditioning the generation on the known fold.  This substantially improves
accuracy (especially for de~novo proteins) but still requires a full AlphaFold forward pass
per conformation per timestep.

\paragraph{ESMFlow-T.}
Rather than using the expensive MSA-based Evoformer, ESMFlow-T uses ESMFold's \citep{lin2023}
protein language model (PLM) embeddings as the sequence representation.  PLM embeddings encode
evolutionary information learned from large-scale protein sequence databases without requiring
per-sequence MSA computation, reducing inference time to $11.2$~s per conformation (Table
in Section~\ref{sec:results}).

\paragraph{Distilled variants.}
\citet{jing2024alphaflow} also train distilled models (AlphaFlow-T$_{\mathrm{dist}}$,
AlphaFlow-T$_{12\mathrm{L,dist}}$) that reduce the number of flow timesteps from 10 to 1
by consistency distillation and decrease model depth, achieving inference times of $3.3$~s
and $1.2$~s respectively, at the cost of somewhat reduced accuracy.

\paragraph{Core limitations of AlphaFlow.}
\begin{enumerate}
  \item \textbf{Dependence on evolutionary information.}  Even AlphaFlow-T requires either
    MSA computation (expensive; unavailable for de~novo proteins) or ESM embeddings
    (trained on natural protein sequences; potentially biased).
  \item \textbf{Inference cost.}  AlphaFlow-T requires a full AlphaFold forward pass at each
    of the 10 flow timesteps: $32.6$~s per conformation.
  \item \textbf{Over-stabilisation.}  AlphaFlow-T systematically under-predicts ensemble
    spread (median RMSF $= 1.17$~Å vs.\ MD value of $1.48$~Å), consistent with the folding
    model's prior bias toward equilibrium structures.
  \item \textbf{Multi-chain proteins.}  AlphaFlow's architecture encodes chain indices as
    features, causing it to fail when parser-modified to accept multi-chain inputs.
\end{enumerate}

%% ───────────────────────────────────────────────────────────────────────────
\subsection{ConfDiff: Force-Guided SE(3) Diffusion}
\label{sec:confdiff}

\citet{wang2024confdiff} proposed ConfDiff, which replaces AlphaFold~2's structure module with
an SE(3) diffusion model while retaining AlphaFold~2's pre-trained Evoformer representations
as conditioning features.  The diffusion model generates conformations conditioned on the
Evoformer embedding of the input sequence and, optionally, a force field-based energy signal.

\paragraph{Key idea.}
Rather than generating structures from sequence directly, ConfDiff generates conformational
perturbations from an equilibrium structure, guided by both the learned distribution and an
energy-based signal.  The energy guidance improves physical realism.

\paragraph{Limitations.}
ConfDiff still requires AlphaFold~2's Evoformer (hence MSA), has inference time of $20.2$~s
per conformation, and underperforms AlphaFlow-T on the ATLAS benchmark in most metrics
(Table in Section~\ref{sec:results}).

%% ───────────────────────────────────────────────────────────────────────────
\subsection{BioEmu: Scalable Equilibrium Ensemble Emulation}
\label{sec:bioemu}

\citet{lewis2025bioemu} took a different approach and trained BioEmu, an AlphaFold-like model,
on a very large heterogeneous dataset comprising MDs of varying lengths, temperatures, and
force fields, supplemented by NMR and crystallographic ensemble data.  BioEmu targets
\emph{general} equilibrium ensembles---including rare alternative folded states and
disordered regions---rather than the specific distribution induced by a fixed MD protocol.

\paragraph{Strengths.}
BioEmu generates diverse ensembles and can sample rare states not accessible in short MD.
Its large training set partially mitigates over-fitting to a single MD protocol.

\paragraph{Limitations for standardised MD emulation.}
Because BioEmu is not trained to reproduce the specific distribution of $300$~ns ATLAS
trajectories, it performs unfavourably on the ATLAS benchmark: RMSF MAE of $1.29$~Å
vs.\ $0.42$~Å for BBFlow (Table in Section~\ref{sec:results}).  This is expected---BioEmu
optimises for a different objective.

\paragraph{Summary.}
Table~\ref{tab:related_summary} summarises key properties of the baselines and BBFlow.

\begin{table}[h]
  \centering
  \caption{Comparison of ensemble generation approaches.  ``MSA'' = requires Multiple
    Sequence Alignment.  ``PLM'' = protein language model embeddings.  ``FT'' = fine-tunes
    pre-trained weights.  ``Struct'' = conditions on equilibrium structure.  ``Multi'' = demonstrated
    on multi-chain proteins.}
  \label{tab:related_summary}
  \begin{tabular}{lcccccc}
    \toprule
    \textbf{Method} & \textbf{MSA} & \textbf{PLM} & \textbf{FT} & \textbf{Struct} &
      \textbf{Multi} & \textbf{Time (s)} \\
    \midrule
    AlphaFlow       & \checkmark & --          & \checkmark & --          & --          & 32.0 \\
    AlphaFlow-T     & \checkmark & --          & \checkmark & \checkmark  & --          & 32.6 \\
    ESMFlow-T       & --         & \checkmark  & \checkmark & \checkmark  & --          & 11.2 \\
    ConfDiff        & \checkmark & --          & \checkmark & --          & --          & 20.2 \\
    BioEmu          & \checkmark & --          & \checkmark & --          & --          &  1.9 \\
    MDGen           & --         & --          & --         & --          & --          &  0.15 \\
    \midrule
    \textbf{BBFlow} & --         & --          & --         & \checkmark  & \checkmark  &  \textbf{0.8} \\
    \bottomrule
  \end{tabular}
\end{table}

%% ───────────────────────────────────────────────────────────────────────────
\subsection{MDGen: Time-Consistent All-Atom Trajectory Generation}
\label{sec:mdgen}

\citet{jing2024mdgen} proposed MDGen, a generative model trained to produce all-atom
trajectories with physically consistent temporal dynamics---not just ensembles of independent
snapshots.  MDGen operates in the discrete token space of molecular structure, modelling
the joint distribution over consecutive frames.

\paragraph{Relation to BBFlow.}
While MDGen and BBFlow both use the ATLAS dataset, their objectives differ: MDGen models
temporal correlations between frames while BBFlow models only the marginal distribution
(independent snapshot generation).  BBFlow achieves substantially higher ensemble accuracy
(RMSF MAE $0.42$ vs.\ $0.81$~Å; DCCM correlation $0.87$ vs.\ $0.54$), while MDGen is
faster ($0.15$ vs.\ $0.8$~s per conformation) due to its token-based architecture.

%% ───────────────────────────────────────────────────────────────────────────
\subsection{Summary of the Gap BBFlow Addresses}
\label{sec:gap}

The landscape of prior work reveals a clear dichotomy:
\begin{itemize}
  \item \textbf{Accurate but slow}: AlphaFlow-T achieves strong RMSF correlation ($0.92$) but
    requires $32.6$~s per conformation and depends on evolutionary information.
  \item \textbf{Fast but inaccurate}: Distilled AlphaFlow-T$_{12\mathrm{L,dist}}$ ($1.2$~s) and
    BioEmu ($1.9$~s) are faster but substantially less accurate on ATLAS metrics.
  \item \textbf{No model handles multi-chain proteins.}
  \item \textbf{All accurate models require evolutionary information}, creating a systematic
    failure mode for de~novo proteins.
\end{itemize}
BBFlow directly addresses all four gaps simultaneously.

%% ═══════════════════════════════════════════════════════════════════════════════
\section{BBFlow: Method}
\label{sec:method}

\subsection{Problem Formulation}
\label{sec:formulation}

We model the conformational ensemble of a protein as the conditional distribution
$p(x | \xeq)$, where $x \in \SE^N$ is a backbone conformation and $\xeq \in \SE^N$ is
the equilibrium (reference) structure.  This framing is natural for MD emulation because:
\begin{enumerate}
  \item MD also requires an initial (equilibrium) structure.
  \item Conditioning on $\xeq$ encodes the fold identity without recourse to the amino-acid
    sequence or evolutionary data.
  \item The same architecture can generate diverse ensembles for structurally distinct proteins
    by presenting different $\xeq$ inputs at inference.
\end{enumerate}

Formally, we seek to learn a conditional flow matching model that samples from $p(x|\xeq)$:
\begin{equation}
  \label{eq:problem}
  x \sim p(x | \xeq) \approx \phi_1(x_0; \xeq), \quad x_0 \sim p_0(\cdot | \xeq),
\end{equation}
where $\phi_1$ is the flow at time $t=1$ (the endpoint of the ODE integration starting
from $x_0$), and $p_0(\cdot|\xeq)$ is a novel conditional prior we introduce below.

\subsection{Conditional Flow Matching for Ensemble Generation}
\label{sec:conditional_fm}

Extending the flow ODE~\eqref{eq:flow_ode} to include the conditioning signal $\xeq$, we
learn a conditional vector field:
\begin{equation}
  \label{eq:cond_vfield}
  v(x, t, \xeq): \mathcal{M} \times [0,1] \times \Meq \to T_x \mathcal{M},
\end{equation}
defining the conditional flow $\phi_t(\cdot | \xeq)$ via:
\begin{equation}
  \label{eq:cond_flow_ode}
  \frac{d}{dt}\phi_t(x | \xeq) = v(\phi_t, t, \xeq), \qquad \phi_0(x|\xeq) = x.
\end{equation}
The conditional training loss is:
\begin{equation}
  \label{eq:cond_fm_loss}
  \Loss_{\mathrm{FM}} = \mathbb{E}\!\left[\norm{v - \vhat(x_t, t, \xeq)}^2_{\SE}\right],
\end{equation}
where the expectation is over:
\begin{equation}
  t \sim U(0,1), \quad (x_1, \xeq) \sim \pdata, \quad x_0 \sim p_0(\cdot|\xeq),
  \quad x_t = \gamma(x_0, x_1, t),
\end{equation}
$v = v_{\SE}(x_t, t | x_1)$ is the ground-truth vector field from Equations~\eqref{eq:vfield_so3}
and~\eqref{eq:vfield_r3}, and the SE(3) norm is:
\begin{equation}
  \label{eq:se3_norm}
  \norm{v}^2_{\SE} = \mathrm{Tr}(v_r v_r^T)/2 + \norm{v_z}^2_2.
\end{equation}
The first term measures distance in $\mathfrak{so}(3)$ (rotational component) and the second
in $\Rthree$ (translational component).

\subsection{Encoding the Equilibrium Structure}
\label{sec:encoding}

The equilibrium structure $\xeq = \{(R^{\mathrm{eq}}_i, z^{\mathrm{eq}}_i)\}_{i=1}^N$ is
encoded as initial edge features of the graph neural network via two complementary representations.

\paragraph{Distance encoding.}
Inspired by the interpretation of evolutionary contact maps as proxies for structural
proximity \citep{lin2023}, we encode pairwise inter-residue distances as binned scalars:
\begin{equation}
  \label{eq:dist_enc}
  s_{ij} = \mathrm{bin}\!\left(\norm{z^{\mathrm{eq}}_i - z^{\mathrm{eq}}_j}_2\right),
\end{equation}
using 22 uniform bins spanning $[0, 20]$~Å.  This yields a 22-dimensional one-hot
edge feature for each residue pair within a cutoff radius.  The distance encoding is
invariant to global rotations and translations, and provides coarse structural information
analogous to a contact map.

\paragraph{Direction encoding.}
To provide geometrically richer, equivariant information, we encode the unit vector from
residue $i$ to residue $j$ expressed in the local frame of residue $i$:
\begin{equation}
  \label{eq:dir_enc}
  e_{ij} = (R^{\mathrm{eq}}_i)^{-1} \cdot \frac{z^{\mathrm{eq}}_i - z^{\mathrm{eq}}_j}
    {\norm{z^{\mathrm{eq}}_i - z^{\mathrm{eq}}_j}_2},
\end{equation}
computed for all pairs with $\norm{z^{\mathrm{eq}}_i - z^{\mathrm{eq}}_j}_2 < 5$~Å.
By transforming the inter-residue direction into the co-rotating local frame of residue $i$,
the direction vector becomes an invariant feature (invariant to global rotation) that can be
directly concatenated with the scalar features $s_{ij}$.  This encoding captures the local
geometry of the backbone neighbourhood in a form analogous to tensor-network equivariant
features \citep{schutt2021}.

\paragraph{Amino-acid identity.}
Amino-acid identities are encoded as 20-dimensional one-hot vectors, projected to 128-dimensional
node embeddings via a linear layer.  This provides information about local backbone degrees of
freedom (e.g.\ proline is rigid, glycine is highly flexible) without providing global sequence
context.  An ablation shows the model remains competitive even without amino-acid identities
(Section~\ref{sec:ablation}).

\subsection{Conditional Prior Distribution}
\label{sec:cond_prior}

Standard flow matching uses an unconditional prior:
$x_0 = (R_0, z_0)$ with $z_0 \sim \mathcal{N}(0, \sigma^2 I)$ and
$R_0 \sim U(\SOthree)$ \citep{yim2023frameflow}.  The flow must then learn to assemble a
coherent protein backbone from completely random noise, which is challenging.

BBFlow introduces a \textbf{conditional prior} $p_0(\cdot | \xeq)$ that biases the initial
state toward the equilibrium structure via geodesic interpolation on $\SE^N$:
\begin{equation}
  \label{eq:cond_prior}
  x_{\mathrm{uncond}} \sim \puncond, \qquad x_0 = \gamma(x_{\mathrm{uncond}}, \xeq, \xi),
\end{equation}
where $\puncond$ is the standard unconditional prior, $\gamma: \SE^N \times \SE^N \times [0,1]
\to \SE^N$ is the geodesic on $\SE^N$ (factoring per-residue into rotation and translation
geodesics), and $\xi \in (0,1)$ is a hyperparameter controlling the proximity to $\xeq$:
\begin{equation}
  \gamma(x_{\mathrm{uncond}}, \xeq, 0) = x_{\mathrm{uncond}}, \qquad
  \gamma(x_{\mathrm{uncond}}, \xeq, 1) = \xeq.
\end{equation}

\paragraph{Interpretation.}
At $\xi = 0$, the conditional prior degenerates to the unconditional prior.  At $\xi = 1$,
all samples start exactly at the equilibrium structure and the flow learns a small perturbation.
At $\xi = 0.2$ (used in all main experiments), the prior sample is geometrically $20\%$ of the
way from random noise toward $\xeq$ along the manifold geodesic---the backbone topology is
already roughly correct, but with substantial structural diversity allowing the flow to generate
different conformations.

\paragraph{Connection to partial denoising.}
This conditional prior is a generalisation of \emph{partial denoising} from diffusion models
\citep{lu2022str2str} to the flow matching framework on $\SE^N$.  Partial denoising starts
the reverse diffusion from an intermediate noisy state rather than pure noise, biasing
generation toward a reference structure.  BBFlow achieves the same effect in the geometrically
correct framework of geodesic interpolation on the manifold.

\paragraph{Effect on training.}
Conditioning the prior on $\xeq$ reduces the burden on the flow vector field: it no longer
needs to assemble the correct topology from scratch, but only needs to generate the appropriate
conformational deviation from the known equilibrium fold.  The ablation in Section~\ref{sec:ablation}
confirms this: removing the conditional prior degrades RMSF MAE from $0.42$ to $0.48$~Å
and pairwise RMSD MAE from $0.77$ to $0.90$~Å.

\subsection{Network Architecture}
\label{sec:architecture}

The vector field network $\vhat_\theta(x_t, t, \xeq)$ is adapted from GAFL
\citep{wagner2024gafl}, itself an extension of FrameDiff \citep{yim2023framediff} and
FrameFlow \citep{yim2023frameflow}.

\paragraph{Inputs.}
\begin{itemize}
  \item \textbf{Node features}: current frame positions $\{(R^t_i, z^t_i)\}_{i=1}^N$,
    flow time $t$ (scalar, broadcast to all nodes), and amino-acid identity embedding.
  \item \textbf{Edge features}: pairwise distances $s_{ij}$ at current time $t$, equilibrium
    distance encoding $s^{\mathrm{eq}}_{ij}$, and equilibrium direction encoding $e_{ij}$.
\end{itemize}
The equilibrium features $s^{\mathrm{eq}}_{ij}$ and $e_{ij}$ are concatenated with the
time-varying pairwise distances as initial edge features, expanding the edge feature
dimensionality by $22 + 3 = 25$ dimensions.

\paragraph{Message-passing blocks.}
Six CFA message-passing blocks consecutively update node and edge representations.  Each block:
\begin{enumerate}
  \item Computes geometric messages from neighbouring residues using CFA (Section~\ref{sec:cfa}).
  \item Aggregates messages to update node features.
  \item Updates frames $T_i$ by applying the predicted incremental rotation and translation.
  \item Updates edge features using updated node features and geometric inter-frame relationships.
\end{enumerate}

\paragraph{Output.}
The final layer outputs predicted target frames $\hat{x}_1 = \{(\hat{R}^1_i, \hat{z}^1_i)\}$,
from which the predicted vector field is recovered using Equations~\eqref{eq:vfield_so3}
and~\eqref{eq:vfield_r3}.

\paragraph{Omission of residue indices.}
In contrast to AlphaFold~2, FrameDiff, FrameFlow, and AlphaFlow, BBFlow does \emph{not} use
residue indices (absolute position in the chain) as node features.  The reasoning is that
the ordering of residues along the chain is already encoded in the equilibrium structure's
geometry: adjacent residues have small inter-$\mathrm{C}_\alpha$ distances, while sequentially
distant residues are typically further apart.  Removing residue indices serves two purposes:
\begin{enumerate}
  \item \textbf{Reduces memorisation}: residue indices provide a unique identifier for each
    position, potentially allowing the model to memorise per-residue behaviour from the
    training set rather than learning generalisable structural patterns.
  \item \textbf{Enables multi-chain generalisation}: residue indices are defined per chain
    and are not meaningfully comparable across chains or between monomeric and multimeric
    proteins.  Removing them allows the model to process multi-chain inputs without
    modification.
\end{enumerate}
The ablation (Section~\ref{sec:ablation}) shows that adding residue indices (variant~c) does
not improve accuracy on monomeric proteins, confirming that geometric conditioning is sufficient.

\paragraph{Parameter count.}
The full BBFlow model has approximately $18.2$~million learnable parameters.  BBFlow-light,
a smaller variant with 3 instead of 6 CFA blocks and reduced feature dimensions (96/48 vs.\
128/64), has $2.5$~million parameters and achieves $200\times$ speedup over AlphaFlow at
modestly reduced accuracy (Section~\ref{sec:ablation}).

\subsection{Training Algorithm}
\label{sec:training_alg}

The complete training procedure is summarised in Algorithm~\ref{alg:training}.

\begin{algorithm}[h]
  \caption{Training of BBFlow}
  \label{alg:training}
  \begin{algorithmic}[1]
    \Require Dataset $\mathcal{D} = \{(\xeq^{(i)}, \{x_1^{(i,m)}\}_{m=1}^M)\}$,
      with equilibrium structure $\xeq \in \SE^N$ and $M$ MD conformations $x_1 \in \SE^N$ per protein
    \Require Model $\hat{x}_\theta$
    \While{training}
      \State $(\xeq, x_1) \sim \mathcal{D}$ \Comment{Sample equilibrium structure and one MD conformation}
      \State $x_0 \sim p_0(\cdot|\xeq)$ \Comment{Sample from conditional prior (Eq.~\ref{eq:cond_prior})}
      \State $t \sim U(0,1)$ \Comment{Sample flow matching time}
      \State $x_t = \gamma(x_0, x_1, t)$ \Comment{Geodesic interpolation on $\SE^N$}
      \State $v = v_{\SE}(x_t, t | x_1)$ \Comment{Ground-truth vector field (Eqs.~\ref{eq:vfield_so3}, \ref{eq:vfield_r3})}
      \State $\hat{x}_1 = \hat{x}_\theta(x_t, t, \xeq)$ \Comment{Model output}
      \State $\hat{v} = v_{\SE}(x_t, t | \hat{x}_1)$ \Comment{Predicted vector field}
      \State $\Loss = \norm{v - \hat{v}}^2_{\SE} + \Loss_{\mathrm{aux}}(x_1, \hat{x}_1)$ \Comment{Flow matching + auxiliary loss}
      \State Update $\theta$ using $\nabla_\theta \Loss$
    \EndWhile
  \end{algorithmic}
\end{algorithm}

\subsection{Inference}
\label{sec:inference}

At inference, given a query protein with equilibrium structure $\xeq$:
\begin{enumerate}
  \item Sample $x_{\mathrm{uncond}} \sim \puncond$ (Gaussian translations, uniform rotations).
  \item Compute $x_0 = \gamma(x_{\mathrm{uncond}}, \xeq, \xi)$ with $\xi = 0.2$.
  \item Integrate the learned flow ODE~\eqref{eq:cond_flow_ode} using 20 uniform timesteps
    (Euler method on $\SE^N$).
  \item The endpoint $x_1 = \phi_1(x_0|\xeq)$ is a sample from the approximated Boltzmann
    distribution $p(x|\xeq)$.
  \item Repeat steps 1--4 independently for each desired conformation (samples are i.i.d.\
    given $\xeq$).
\end{enumerate}
The use of only 20 ODE integration steps (vs.\ $\sim\!50$--$1000$ for diffusion models)
is enabled by the nearly straight geodesic trajectories characteristic of flow matching.

%% ═══════════════════════════════════════════════════════════════════════════════
\section{Experimental Setup}
\label{sec:experiments}

\subsection{Dataset: ATLAS}
\label{sec:atlas}

The ATLAS database \citep{vandermeersche2024atlas} provides standardised all-atom MD
trajectories for 1390 structurally diverse, monomeric proteins from the Protein Data Bank.
Each protein has three independent $100$-ns trajectories simulated at $310$~K using the
CHARMM36m force field \citep{huang2017} and explicit TIP3P water \citep{jorgensen1983}
in GROMACS \citep{abraham2015gromacs}.  The standardised simulation protocol ensures that
ensemble metrics are comparable across proteins and models.

Following \citet{jing2024alphaflow}, ATLAS is split into 1265 training, 39 validation,
and 82 test proteins.  For evaluation, 250 independent conformations are generated per
protein and compared to the MD trajectories.

\subsection{De Novo Protein Dataset}
\label{sec:denovo_data}

To evaluate generalisation to proteins without natural evolutionary history, we generated
a dataset of 50 de~novo proteins using RFdiffusion \citep{watson2023} and FrameFlow \citep{yim2023frameflow},
each subjected to three $100$-ns MD simulations following the ATLAS protocol.  Equilibrium
structures were obtained using ESMFold \citep{lin2023} (since de~novo sequences have no
natural MSA).  This dataset provides a challenging test because AlphaFlow-style models cannot
access meaningful evolutionary information for these proteins.

\subsection{Multi-Chain Systems}
\label{sec:multichain_data}

Five well-characterised multi-chain systems were simulated:
\begin{itemize}
  \item Barnase-barstar complex (PDB: 1BGS), homodimer with two chains.
  \item Small homotrimeric foldon (PDB: 1RFO).
  \item Homodimeric fructokinase in apo state (PDB: 5EY7).
  \item Heterodimeric nanobody-SARS-CoV-2 RBD complex (PDB: 7KGK).
  \item Heterodimeric nanobody-TNFRSF17 complex (PDB: 8HXR).
\end{itemize}
Systems span 80 to 590 total residues.  BBFlow was trained only on monomeric ATLAS proteins
and applied to these complexes without modification.

\subsection{Evaluation Metrics}
\label{sec:metrics}

All metrics are computed on C$_\alpha$ atoms after superimposing all generated conformations
to the equilibrium structure.  Bootstrap resampling of MD conformations ($100$ replicates)
provides uncertainty estimates for MD-derived statistics.

\paragraph{Root Mean Square Fluctuation (RMSF).}
For each residue $i$, RMSF measures the time-averaged deviation of C$_\alpha$ from its mean
position across the ensemble:
\begin{equation}
  \mathrm{RMSF}_i = \sqrt{\frac{1}{M}\sum_{m=1}^M \norm{z_i^{(m)} - \bar{z}_i}^2},
  \quad \bar{z}_i = \frac{1}{M}\sum_{m=1}^M z_i^{(m)}.
\end{equation}
We report the \textbf{Pearson correlation} $r$ between the MD and generated RMSF profiles
(shape agreement), the \textbf{mean absolute error} (MAE) of per-residue RMSF (amplitude
agreement), and the \textbf{median RMSF} across all residues (over-/under-stabilisation
detection).

\paragraph{Pairwise RMSD.}
The average C$_\alpha$ RMSD between all pairs of conformations quantifies ensemble spread
without requiring a reference state:
\begin{equation}
  \mathrm{pwRMSD} = \frac{1}{M^2}\sum_{m,m'=1}^M \mathrm{RMSD}(x^{(m)}, x^{(m')}).
\end{equation}
We report the MAE of pairwise RMSD between generated and MD ensembles.

\paragraph{Dynamical Cross-Correlation Matrix (DCCM).}
The DCCM encodes correlated motions between pairs of residues:
\begin{equation}
  \mathrm{DCCM}_{ij} = \frac{(z_i - \langle z_i\rangle)\cdot(z_j - \langle z_j\rangle)}
    {\sqrt{\langle(z_i - \langle z_i\rangle)^2\rangle}\sqrt{\langle(z_j - \langle z_j\rangle)^2\rangle}}.
\end{equation}
Values near $+1$ indicate correlated motion (hinge regions), near $-1$ anti-correlated (breathing
motion), and near $0$ uncorrelated.  We report the Pearson correlation $r$ between flattened
DCCM matrices from MD and generated ensembles.

\paragraph{Principal Component Analysis (PCA).}
We compute the first two principal components of the MD C$_\alpha$ trajectory, project both
MD and generated conformations onto these components, and report the \textbf{Wasserstein-2
distance} $W_2$ between the two 2D distributions.  This metric captures whether the model
explores the correct regions of conformational space.

\paragraph{Transient contact accuracy.}
A transient contact is a residue pair that is separated ($> 8$~Å C$_\alpha$ distance) in
the equilibrium structure but in contact ($< 8$~Å) in at least $10\%$ of ensemble conformations.
We report the \textbf{Jaccard similarity} $J_{\mathrm{tr}}$ between the sets of transient
contacts identified in the MD and generated ensembles.

%% ═══════════════════════════════════════════════════════════════════════════════
\section{Results}
\label{sec:results}

\subsection{ATLAS Benchmark}
\label{sec:atlas_results}

Table~\ref{tab:main_results} reports performance on the ATLAS test set.  All models generate
250 conformations per protein; inference time is measured per conformation for the
302-residue protein 7c45A on an NVIDIA A100-80GB GPU.

\begin{table}[h]
  \centering
  \caption{Performance on the ATLAS test set (82 proteins).  Median over all proteins.
    MD reference RMSF median: $1.48$~Å.  $\uparrow$: higher is better.
    $\downarrow$: lower is better.  Bold: best; underline: second best.
    Errors in parentheses when above rounding precision.
    BioEmu$^*$: not trained on ATLAS; not directly comparable.}
  \label{tab:main_results}
  \small
  \begin{tabular}{lccccccc}
    \toprule
    & \multicolumn{2}{c}{\textbf{RMSF}} & \textbf{Pw-RMSD} & \textbf{DCCM}
      & \textbf{PCA} & $\bm{J_{\mathrm{tr}}}$ & \textbf{Time} \\
    \cmidrule(lr){2-3}
    \textbf{Method} & $r\,(\uparrow)$ & MAE$\,(\downarrow)$ & MAE$\,(\downarrow)$
      & $r\,(\uparrow)$ & $W_2\,(\downarrow)$ & $\%\,(\uparrow)$ & [s]$\,(\downarrow)$ \\
    \midrule
    BioEmu$^*$            & 0.83 & 1.29         & 2.84         & 0.80 & 1.65         & 36 & 1.9 \\
    AlphaFlow             & 0.86 & 0.59         & 1.35         & 0.86 & 1.47         & 41 & 32.0 \\
    ConfDiff              & 0.88 & 0.62         & 1.45         & 0.86 & 1.41         & 39 & 20.2 \\
    AlphaFlow-T           & 0.92 & 0.41         & 0.91         & 0.89 & 1.28         & 47 & 32.6 \\
    ESMFlow-T             & 0.92 & 0.52         & 1.22         & 0.89 & 1.48         & 47 & 11.2 \\
    AlphaFlow-T$_{\mathrm{dist}}$ & 0.92 & 0.68 & 1.41        & 0.88 & 1.43         & 42 & 3.3 \\
    AlphaFlow-T$_{12\mathrm{L,dist}}$ & 0.90 & 0.85 & 1.80    & 0.87 & 1.60         & 24 & 1.2 \\
    \midrule
    \textbf{BBFlow}       & \underline{0.90} & \underline{0.42} & \underline{0.77}
      & \underline{0.87} & \underline{1.33} & 29 & \textbf{0.8} \\
    \bottomrule
  \end{tabular}
\end{table}

\paragraph{Key observations.}
\begin{enumerate}
  \item \textbf{Speed vs.\ accuracy Pareto front.}
    BBFlow occupies a unique position: it is $41\times$ faster than AlphaFlow-T while
    achieving comparable or superior accuracy on five of six metrics.

  \item \textbf{Over-stabilisation in AlphaFlow-T.}
    AlphaFlow-T has a median RMSF of $1.17$~Å vs.\ the MD reference of $1.48$~Å---a
    systematic under-prediction of ensemble spread.  BBFlow's median RMSF of $1.49$~Å
    matches MD almost exactly.  This over-stabilisation in AlphaFlow-T is attributed to the
    folding model's strong prior toward equilibrium structures.

  \item \textbf{Pairwise RMSD.}
    BBFlow achieves the best pairwise RMSD MAE ($0.77$~Å), demonstrating superior capture of
    the overall magnitude of conformational changes compared to all baselines including
    AlphaFlow-T ($0.91$~Å).

  \item \textbf{RMSF correlation vs.\ MAE.}
    AlphaFlow-T achieves RMSF correlation $r = 0.92$ vs.\ BBFlow's $0.90$, indicating
    slightly better capture of which residues are flexible.  However, BBFlow's RMSF MAE
    ($0.42$~Å) matches AlphaFlow-T ($0.41$~Å), indicating that the amplitudes are equally
    accurate.

  \item \textbf{Transient contacts.}
    BBFlow underperforms MSA-based models on $J_{\mathrm{tr}}$ ($29\%$ vs.\ AlphaFlow-T's
    $47\%$).  This is expected: transient contacts are rare events that require the model to
    predict large-amplitude excursions from equilibrium.  MSA-derived co-evolutionary signals
    carry information about alternative contact patterns that BBFlow's geometry-only conditioning
    does not.

  \item \textbf{Scaling with protein length.}
    BBFlow's advantage over baselines grows with protein length: for large proteins ($> 350$
    residues), BBFlow's RMSF and pairwise RMSD metrics are substantially better, consistent
    with the folding-model baselines over-stabilising large flexible proteins.
\end{enumerate}

\subsection{De Novo Proteins}
\label{sec:denovo_results}

Table~\ref{tab:denovo_results} reports performance on the de~novo protein dataset.

\begin{table}[h]
  \centering
  \caption{Performance on de~novo proteins (50 systems).  Evaluation settings as in
    Table~\ref{tab:main_results}.  MD reference RMSF median: $0.91$~Å.}
  \label{tab:denovo_results}
  \small
  \begin{tabular}{lccccccc}
    \toprule
    & \multicolumn{2}{c}{\textbf{RMSF}} & \textbf{Pw-RMSD} & \textbf{DCCM}
      & \textbf{PCA} & $\bm{J_{\mathrm{tr}}}$ & \textbf{Time} \\
    \cmidrule(lr){2-3}
    \textbf{Method} & $r\,(\uparrow)$ & MAE$\,(\downarrow)$ & MAE$\,(\downarrow)$
      & $r\,(\uparrow)$ & $W_2\,(\downarrow)$ & $\%\,(\uparrow)$ & [s]$\,(\downarrow)$ \\
    \midrule
    BioEmu$^*$        & 0.60 & 4.24 & 8.29 & 0.64 & 1.53 & 23 & 1.9 \\
    AlphaFlow         & 0.47 & 4.76 & 7.40 & 0.58 & 1.64 & 17 & 32.0 \\
    ConfDiff          & 0.62 & 3.82 & 7.26 & 0.65 & 1.72 & 15 & 20.2 \\
    AlphaFlow-T       & \underline{0.89} & \textbf{0.25} & 0.38 & \underline{0.85} & \textbf{0.66} & \textbf{55} & 32.6 \\
    ESMFlow-T         & \underline{0.89} & \underline{0.28} & \underline{0.43} & \textbf{0.86} & \underline{0.63} & \textbf{55} & 11.2 \\
    AlphaFlow-T$_{\mathrm{dist}}$ & 0.88 & 0.46 & 0.77 & 0.84 & 0.69 & 51 & 3.3 \\
    AlphaFlow-T$_{12\mathrm{L,dist}}$ & 0.87 & 0.58 & 0.97 & 0.83 & 0.75 & 38 & 1.2 \\
    \midrule
    \textbf{BBFlow}   & 0.84 & 0.26 & \textbf{0.32} & 0.83 & 0.67 & 32 & \textbf{0.8} \\
    \bottomrule
  \end{tabular}
\end{table}

\paragraph{Key observations.}
\begin{enumerate}
  \item \textbf{Catastrophic failure of evolutionary-information-dependent models.}
    AlphaFlow (no templates) degrades from RMSF correlation $0.86$ on natural proteins to
    $0.47$ on de~novo proteins---essentially random.  BioEmu similarly degrades from $0.83$
    to $0.60$.  These models rely on MSA or PLM embeddings for which de~novo proteins have
    no meaningful data, causing the conditioning signal to be uninformative or misleading.

  \item \textbf{BBFlow maintains competitive performance.}
    BBFlow's RMSF correlation drops only from $0.90$ to $0.84$ ($\Delta = 0.06$) compared to
    AlphaFlow's drop of $0.39$.  Its RMSF MAE of $0.26$~Å is near-identical to
    AlphaFlow-T's best of $0.25$~Å, despite $41\times$ faster inference.

  \item \textbf{Pairwise RMSD: BBFlow is best.}
    BBFlow achieves the lowest pairwise RMSD MAE ($0.32$~Å) among all models, indicating
    that it best captures the true magnitude of conformational diversity in de~novo proteins.

  \item \textbf{AlphaFlow-T retains advantage with templates.}
    Template-conditioned AlphaFlow variants still perform well on de~novo proteins because
    the equilibrium structure is provided as a template, bypassing the need for sequence-based
    fold prediction.  However, their inference cost ($32.6$~s) remains prohibitive for
    screening applications.
\end{enumerate}

\subsection{Multi-Chain Proteins}
\label{sec:multichain_results}

Table~\ref{tab:multichain_results} reports BBFlow's performance on multi-chain systems.
No other baseline was able to generate physically valid conformations for these systems.

\begin{table}[h]
  \centering
  \caption{BBFlow performance on five multi-chain proteins.
    MD reference RMSF median: $1.20$~Å.  Settings as in Table~\ref{tab:main_results}.}
  \label{tab:multichain_results}
  \begin{tabular}{lccccc}
    \toprule
    & \multicolumn{2}{c}{\textbf{RMSF}} & \textbf{Pw-RMSD} & \textbf{DCCM} & \textbf{PCA} \\
    \cmidrule(lr){2-3}
    \textbf{Method} & $r\,(\uparrow)$ & MAE$\,(\downarrow)$ & MAE$\,(\downarrow)$
      & $r\,(\uparrow)$ & $W_2\,(\downarrow)$ \\
    \midrule
    BBFlow & 0.82 & 0.31 & 0.41 & 0.85 & 0.71 \\
    \bottomrule
  \end{tabular}
\end{table}

\paragraph{Key observations.}
BBFlow achieves RMSF correlation $0.82$ and DCCM correlation $0.85$ on multi-chain systems,
despite having been trained exclusively on monomeric proteins.  The DCCM analysis is particularly
compelling: BBFlow correctly captures both \emph{intra-chain} correlated motions and
\emph{inter-chain} allosteric communication between subunits (e.g.\ correlated hinge motion
in the barnase-barstar complex), which are physically meaningful but absent from training data.

AlphaFlow, when parser-modified to accept multi-chain inputs, generates unphysical states
with broken chain connectivity---a consequence of the chain index features encoding an
assumption of single-chain topology.

\subsection{Sequence-to-Ensemble Pipeline}
\label{sec:seq_to_ensemble}

When only a sequence is available (no equilibrium structure), BBFlow can be combined with
a folding model in a two-stage pipeline: AlphaFold~2 predicts $\xeq$, which is then passed
to BBFlow.  Table~\ref{tab:seq2ens} shows that this pipeline outperforms AlphaFlow
(sequence-to-structure model) in most metrics while being $\sim\!30\times$ faster, because
AlphaFold~2 runs only once (not once per conformation per timestep).

\begin{table}[h]
  \centering
  \caption{Sequence-to-ensemble pipeline comparison on ATLAS test set.}
  \label{tab:seq2ens}
  \small
  \begin{tabular}{lcccccc}
    \toprule
    \textbf{Method} & RMSF $r$ & RMSF MAE & Pw-RMSD MAE & DCCM $r$ & PCA $W_2$ & Time [s] \\
    \midrule
    AlphaFlow           & 0.86 & 0.59 & 1.35 & 0.86 & 1.47 & 32.0 \\
    AlphaFold2 + BBFlow & 0.87 & 0.52 & 1.07 & 0.85 & 1.47 & $0.3 + 0.8$ \\
    \bottomrule
  \end{tabular}
\end{table}

%% ═══════════════════════════════════════════════════════════════════════════════
\section{Ablation Study}
\label{sec:ablation}

Table~\ref{tab:ablation} quantifies the individual contribution of each BBFlow component
by systematically removing or altering key features and evaluating on the ATLAS test set.

\begin{table}[h]
  \centering
  \caption{Ablation study.  Each variant removes or replaces one component of BBFlow.
    ``Cond.\ prior'': conditional prior (Section~\ref{sec:cond_prior}).
    ``Dist.'': distance encoding (Eq.~\ref{eq:dist_enc}).
    ``Dir.'': direction encoding (Eq.~\ref{eq:dir_enc}).
    ``AA'': amino-acid identity.
    ``Idx'': residue index encoding.
    \checkmark: present; --: absent.}
  \label{tab:ablation}
  \begin{tabular}{lccccccc}
    \toprule
    \textbf{Variant} & \textbf{Cond.\ prior} & \textbf{Dist.} & \textbf{Dir.} & \textbf{AA}
      & \textbf{Idx} & RMSF MAE & Pw-RMSD MAE \\
    \midrule
    BBFlow  & \checkmark & \checkmark & \checkmark & \checkmark & --           & \textbf{0.42} & \textbf{0.77} \\
    (a) no direction & -- & \checkmark & --         & \checkmark & \checkmark   & 0.52 & 1.15 \\
    (b) no cond.\ prior & -- & \checkmark & \checkmark & \checkmark & \checkmark & 0.48 & 0.90 \\
    (c) with index  & \checkmark & \checkmark & \checkmark & \checkmark & \checkmark & 0.42 & 0.82 \\
    (d) no AA       & \checkmark & \checkmark & --         & --         & \checkmark & 0.54 & 0.93 \\
    (e) no distance & \checkmark & --         & --         & \checkmark & \checkmark & 5.88 & 7.08 \\
    \bottomrule
  \end{tabular}
\end{table}

\paragraph{Distance encoding (e): essential.}
Removing the distance encoding of the equilibrium structure (variant~e) is catastrophic: RMSF
MAE increases from $0.42$ to $5.88$~Å and pairwise RMSD MAE from $0.77$ to $7.08$~Å---worse
than random generation.  The distance encoding is the primary signal that communicates
\emph{which} fold to generate.  Without it, the model has no information about the target
topology at the edge-feature level.

\paragraph{Direction encoding (a): important.}
Removing the direction encoding (variant~a) degrades RMSF MAE from $0.42$ to $0.52$~Å and
pairwise RMSD MAE from $0.77$ to $1.15$~Å.  The direction encoding provides geometrically
richer, equivariant information about local backbone topology beyond what scalar distances
alone can express, enabling better discrimination of local structural features.

\paragraph{Conditional prior (b): beneficial.}
Removing the conditional prior (variant~b, reverting to the unconditional prior) degrades
RMSF MAE from $0.42$ to $0.48$~Å and pairwise RMSD MAE from $0.77$ to $0.90$~Å.  While
the improvement may appear modest in isolation, the conditional prior also substantially
improves convergence speed during training (fewer epochs to reach equivalent accuracy).

\paragraph{Residue index (c): neutral for monomers, harmful for multimers.}
Adding residue indices (variant~c) does not improve accuracy on the monomeric ATLAS benchmark
(RMSF MAE identical; pairwise RMSD MAE slightly worse at $0.82$).  This confirms that
geometric conditioning is sufficient for monomeric proteins.  However, residue indices cause
failure on multi-chain proteins, as described in Section~\ref{sec:multichain_results}.

\paragraph{Amino-acid identity (d): helpful but not essential.}
Removing amino-acid identities (variant~d) degrades RMSF MAE from $0.42$ to $0.54$~Å.
The model remains competitive---outperforming AlphaFlow ($0.59$~Å) and all distilled models---
even without any sequence information.  This demonstrates that backbone geometry alone is
sufficient for high-quality ensemble generation and opens the door to purely structure-based
applications.

\paragraph{Hyperparameter $\xi$ for the conditional prior.}
A systematic study of the interpolation parameter $\xi$ (Table~\ref{tab:xi_ablation})
reveals a trade-off: small $\xi$ produces diverse but inaccurate ensembles (the flow must
learn a large transformation from noise to ensemble); large $\xi$ constrains the flow toward
the equilibrium structure, producing accurate but under-diverse ensembles.  The optimum
is $\xi = 0.2$, which balances accuracy and diversity.

\begin{table}[h]
  \centering
  \caption{Inference-time ablation of $\xi$ (model trained with $\xi = 0.2$).
    Units as in Table~\ref{tab:main_results}.}
  \label{tab:xi_ablation}
  \begin{tabular}{lcccccc}
    \toprule
    $\xi$ & RMSF $r$ & RMSF MAE & RMSF Med. & Pw-RMSD MAE & DCCM $r$ & PCA $W_2$ \\
    \midrule
    0.01 & 0.70 & 7.61 & 11.16 & 10.83 & 0.73 & 2.85 \\
    0.05 & 0.81 & 3.89 &  5.43 &  5.38 & 0.78 & 2.06 \\
    0.10 & 0.88 & 1.32 &  2.62 &  2.21 & 0.84 & 1.51 \\
    \textbf{0.20} & \textbf{0.90} & \textbf{0.42} & \textbf{1.49} & \textbf{0.77}
      & \textbf{0.87} & \textbf{1.33} \\
    0.30 & 0.89 & 0.47 &  1.37 &  1.02 & 0.86 & 1.69 \\
    0.40 & 0.86 & 0.64 &  1.65 &  1.31 & 0.80 & 2.12 \\
    0.50 & 0.79 & 0.80 &  1.82 &  1.50 & 0.74 & 3.00 \\
    0.60 & 0.73 & 0.82 &  1.70 &  1.65 & 0.70 & 4.51 \\
    \bottomrule
  \end{tabular}
\end{table}

%% ═══════════════════════════════════════════════════════════════════════════════
\section{Discussion}
\label{sec:discussion}

\subsection{What Makes BBFlow Novel?}
\label{sec:novelty}

BBFlow contributes three conceptually distinct innovations to the field of protein conformational
ensemble generation:

\paragraph{1.\ Decoupling ensemble generation from sequence modelling.}
All prior transferable ensemble generation models (AlphaFlow, ESMFlow, ConfDiff, BioEmu)
inherit the architecture of protein folding models, fundamentally framing ensemble generation
as a sequence-to-structure problem.  This coupling is not necessary: the fold of a protein
is already captured in its equilibrium structure, which can be used directly as a geometric
conditioning signal.  BBFlow demonstrates that the sequence-based machinery (MSA computation,
Evoformer, language model embeddings) is not a prerequisite for high-quality ensemble generation,
and that the resulting model is faster, more transferable, and less biased.

\paragraph{2.\ Conditional prior via geodesic interpolation.}
The concept of a geometry-conditioned prior distribution in flow matching is novel.  Prior
work either used unconditional priors (GAFL, FrameFlow) or non-Gaussian unconditional priors
(Chroma \citep{ingraham2023}).  BBFlow's conditional prior generalises partial denoising from
diffusion models to the manifold flow matching framework, providing a principled and geometrically
correct way to bias the generation toward a known reference structure.

\paragraph{3.\ First ensemble model for multi-chain proteins.}
The observation that removing residue-index features enables zero-shot generalisation from
monomers to multimers is non-obvious and has significant practical implications.  Multi-chain
proteins---heterodimers, trimers, antibody-antigen complexes---are central to drug discovery
and protein engineering, yet no prior ensemble generation model had addressed this setting.

\subsection{Limitations}
\label{sec:limitations}

\paragraph{MD distribution emulation.}
BBFlow is fundamentally an emulator of the MD-derived conformational distribution.  It
reproduces the ensemble that would be obtained from three $100$-ns CHARMM36m GROMACS
simulations at $310$~K.  It cannot predict:
\begin{itemize}
  \item \emph{Rare alternative folded states} not sampled in short MD.
  \item \emph{Millisecond-timescale transitions} (e.g.\ helix-coil transitions, large
    domain movements).
  \item \emph{Protein unfolding} or denatured ensembles.
  \item \emph{Ligand binding-induced conformational change} (requires training data with ligand-bound
    MD trajectories).
\end{itemize}
For these applications, models trained on longer or more diverse MD data (BioEmu), or
physics-based enhanced sampling methods, remain more appropriate.

\paragraph{Transient contacts.}
The systematic underperformance on transient contact accuracy ($J_{\mathrm{tr}}$) suggests
that rare large-amplitude excursions from equilibrium, which create new inter-residue contacts,
are better predicted by models with access to evolutionary co-variation patterns.  This is
consistent with the interpretation that co-evolutionary signals encode information about
alternative structural states that proteins can adopt---information not present in a single
equilibrium structure.

\paragraph{Backbone-only generation.}
BBFlow generates only backbone (C$_\alpha$) coordinates.  Side-chain conformational ensembles
and protein--ligand interactions require additional modelling and are subjects for future work.

\paragraph{Input structure dependency.}
BBFlow requires an equilibrium structure as input.  This is not a fundamental limitation for
MD emulation (MD also requires an initial structure), and BBFlow can be combined with a
folding model for sequence-only inputs.  However, it limits direct application to proteins
for which no structural model is available and for which folding models fail.

\subsection{Robustness to Input Quality}
\label{sec:robustness}

An important practical question is whether BBFlow's performance degrades gracefully when
the input equilibrium structure is imperfect (e.g.\ obtained from a folding model prediction
rather than X-ray crystallography).  The authors evaluated BBFlow with Gaussian noise added
to equilibrium C$_\alpha$ coordinates (standard deviations of $0.1$~Å and $0.2$~Å) and with
random MD conformations substituted as the equilibrium structure.  In all cases, performance
degraded only modestly (RMSF MAE increasing from $0.42$ to $0.47$--$0.48$~Å), demonstrating
that BBFlow is robust to small inaccuracies in the input structure.

\subsection{Practical Relevance}
\label{sec:practical}

The combination of accuracy, speed, and transferability positions BBFlow as a practically
useful tool for:
\begin{itemize}
  \item \textbf{Drug discovery}: rapid screening of protein--ligand binding-site flexibility
    to identify cryptic pockets or assess induced-fit binding.
  \item \textbf{Antibody engineering}: evaluation of multi-chain complex flexibility for
    antibody design (nanobody-antigen systems are demonstrated in this work).
  \item \textbf{De~novo protein design}: integration into design pipelines (e.g.\
    RFdiffusion + ProteinMPNN + BBFlow) to screen designed sequences for desired
    dynamic properties---a capability that does not exist in current design pipelines.
  \item \textbf{Large-scale annotation}: at $0.8$~s per conformation, generating $250$
    conformations for all $\sim 200{,}000$ proteins in the AlphaFold database would require
    approximately $10$~GPU-days---a feasible task.
\end{itemize}

%% ═══════════════════════════════════════════════════════════════════════════════
\section{Conclusion}
\label{sec:conclusion}

We have presented a comprehensive technical exposition of BBFlow, a flow matching model for
generating MD-like protein backbone conformational ensembles conditioned on equilibrium backbone
geometry.  The model's key innovations---conditional prior via geodesic interpolation, geometric
encoding of the equilibrium structure as invariant pairwise distances and equivariant local
directions, and exclusion of residue-index features---collectively enable state-of-the-art
performance in the speed/accuracy trade-off for conformational ensemble generation.

BBFlow achieves $41\times$ speedup over the best accuracy-comparable baseline (AlphaFlow-T)
while generating more physically accurate ensembles (better pairwise RMSD, better median
RMSF calibration), demonstrating that large pre-trained folding models and evolutionary sequence
information are not necessary prerequisites for high-quality MD emulation.  The model
generalises to de~novo proteins without degradation, and is the first ensemble generation
model demonstrated on multi-chain protein complexes.

The two core ideas introduced in this work---structure conditioning as a replacement for
sequence-level evolutionary information, and conditional priors in Riemannian flow matching---are
broadly applicable beyond protein ensemble generation and may inform future work in
structural biology, drug discovery, and generative modelling on manifolds.

Code and pre-trained weights are publicly available at
\url{https://github.com/graeter-group/bbflow}.

%% ═══════════════════════════════════════════════════════════════════════════════
\begin{thebibliography}{99}

\bibitem{abraham2015gromacs}
M.\ J.\ Abraham, T.\ Murtola, R.\ Schulz, S.\ P\'{a}ll, J.\ C.\ Smith, B.\ Hess, and E.\ Lindahl.
\newblock {GROMACS}: High performance molecular simulations through multi-level parallelism from laptops to supercomputers.
\newblock \emph{SoftwareX}, 1--2:19--25, 2015.

\bibitem{adcock2006}
S.\ A.\ Adcock and J.\ A.\ McCammon.
\newblock Molecular dynamics: Survey of methods for simulating the activity of proteins.
\newblock \emph{Chemical Reviews}, 106(5):1589--1615, 2006.

\bibitem{baek2021}
M.\ Baek et~al.
\newblock Accurate prediction of protein structures and interactions using a three-track neural network.
\newblock \emph{Science}, 373(6557):871--876, 2021.

\bibitem{benkovic2008}
S.\ J.\ Benkovic, G.\ G.\ Hammes, and S.\ Hammes-Schiffer.
\newblock Free-energy landscape of enzyme catalysis.
\newblock \emph{Biochemistry}, 47(11):3317--3321, 2008.

\bibitem{bernetti2024}
M.\ Bernetti et~al.
\newblock Probing allosteric communication with combined molecular dynamics simulations and network analysis.
\newblock \emph{Current Opinion in Structural Biology}, 86:1--10, 2024.

\bibitem{bose2024}
J.\ Bose et~al.
\newblock {SE(3)}-stochastic flow matching for protein backbone generation.
\newblock In \emph{ICLR 2024}, 2024.

\bibitem{case1994}
D.\ A.\ Case.
\newblock Normal mode analysis of protein dynamics.
\newblock \emph{Current Opinion in Structural Biology}, 4(2):285--290, 1994.

\bibitem{chen2024}
R.\ T.\ Q.\ Chen and Y.\ Lipman.
\newblock Flow matching on general geometries.
\newblock In \emph{ICLR 2024}, 2024.

\bibitem{fuchs2020}
F.\ B.\ Fuchs, D.\ E.\ Worrall, V.\ Fischer, and M.\ Welling.
\newblock {SE(3)}-transformers: 3D roto-translation equivariant attention networks.
\newblock In \emph{NeurIPS 2020}, 2020.

\bibitem{ingraham2023}
J.\ B.\ Ingraham et~al.
\newblock Illuminating protein space with a programmable generative model.
\newblock \emph{Nature}, pages 1--9, 2023.

\bibitem{jing2024alphaflow}
B.\ Jing, B.\ Berger, and T.\ Jaakkola.
\newblock {AlphaFold} meets flow matching for generating protein ensembles.
\newblock In \emph{ICML 2024}, 2024.

\bibitem{jing2024mdgen}
B.\ Jing, H.\ Stark, T.\ Jaakkola, and B.\ Berger.
\newblock Generative modeling of molecular dynamics trajectories.
\newblock In \emph{NeurIPS 2024}, 2024.

\bibitem{jorgensen1983}
W.\ L.\ Jorgensen et~al.
\newblock Comparison of simple potential functions for simulating liquid water.
\newblock \emph{The Journal of Chemical Physics}, 79(2):926--935, 1983.

\bibitem{jumper2021}
J.\ Jumper et~al.
\newblock Highly accurate protein structure prediction with {AlphaFold}.
\newblock \emph{Nature}, 596(7873):583--589, 2021.

\bibitem{klein2023}
L.\ Klein, A.\ Kr\"{a}mer, and F.\ No\'{e}.
\newblock Equivariant flow matching.
\newblock In \emph{NeurIPS 2023}, 2023.

\bibitem{lewis2025bioemu}
S.\ Lewis et~al.
\newblock Scalable emulation of protein equilibrium ensembles with generative deep learning.
\newblock \emph{Science}, 389(6761):eadv9817, 2025.

\bibitem{lin2023}
Z.\ Lin et~al.
\newblock Evolutionary-scale prediction of atomic-level protein structure with a language model.
\newblock \emph{Science}, 379(6637):1123--1130, 2023.

\bibitem{lipman2023}
Y.\ Lipman, R.\ T.\ Q.\ Chen, H.\ Ben-Hamu, M.\ Nickel, and M.\ Le.
\newblock Flow matching for generative modeling.
\newblock In \emph{ICLR 2023}, 2023.

\bibitem{lu2022str2str}
J.\ Lu et~al.
\newblock {Str2Str}: A score-based framework for zero-shot protein conformation sampling.
\newblock In \emph{ICLR 2024}, 2024.

\bibitem{noe2019}
F.\ No\'{e}, S.\ Olsson, J.\ K\"{o}hler, and H.\ Wu.
\newblock Boltzmann generators: Sampling equilibrium states of many-body systems with deep learning.
\newblock \emph{Science}, 365(6457):eaaw1147, 2019.

\bibitem{puny2022}
O.\ Puny et~al.
\newblock Frame averaging for invariant and equivariant network design.
\newblock In \emph{ICLR 2022}, 2022.

\bibitem{satorras2021}
V.\ G.\ Satorras, E.\ Hoogeboom, and M.\ Welling.
\newblock E(n) equivariant graph neural networks.
\newblock In \emph{ICML 2021}, 2021.

\bibitem{schutt2021}
K.\ T.\ Sch\"{u}tt, O.\ T.\ Unke, and M.\ Gastegger.
\newblock Equivariant message passing for the prediction of tensorial properties and molecular spectra.
\newblock In \emph{ICML 2021}, 2021.

\bibitem{song2021}
Y.\ Song et~al.
\newblock Score-based generative modeling through stochastic differential equations.
\newblock In \emph{ICLR 2021}, 2021.

\bibitem{thomas2018}
N.\ Thomas et~al.
\newblock Tensor field networks: Rotation- and translation-equivariant neural networks for 3D point clouds.
\newblock \emph{arXiv:1802.08219}, 2018.

\bibitem{vandermeersche2024atlas}
Y.\ Vander Meersche, G.\ Cretin, A.\ Gheeraert, J.-C.\ Gelly, and T.\ Galochkina.
\newblock {ATLAS}: Protein flexibility description from atomistic molecular dynamics simulations.
\newblock \emph{Nucleic Acids Research}, 52(D1):D384--D392, 2024.

\bibitem{vaswani2017}
A.\ Vaswani et~al.
\newblock Attention is all you need.
\newblock In \emph{NeurIPS 2017}, 2017.

\bibitem{wagner2024gafl}
S.\ Wagner, L.\ Seute, V.\ Viliuga, N.\ Wolf, F.\ Gr\"{a}ter, and J.\ Stuehmer.
\newblock Generating highly designable proteins with geometric algebra flow matching.
\newblock In \emph{NeurIPS 2024}, 2024.

\bibitem{wang2024confdiff}
Y.\ Wang et~al.
\newblock Protein conformation generation via force-guided {SE(3)} diffusion models.
\newblock In \emph{ICML 2024}, 2024.

\bibitem{watson2023}
J.\ L.\ Watson et~al.
\newblock De novo design of protein structure and function with {RFdiffusion}.
\newblock \emph{Nature}, 620:1089--1100, 2023.

\bibitem{wayment2024}
H.\ K.\ Wayment-Steele et~al.
\newblock Predicting multiple conformations via sequence clustering and {AlphaFold2}.
\newblock \emph{Nature}, 625:832--839, 2024.

\bibitem{wolf2025bbflow}
N.\ Wolf, L.\ Seute, V.\ Viliuga, S.\ Wagner, J.\ St\"{u}hmer, and F.\ Gr\"{a}ter.
\newblock Learning conformational ensembles of proteins based on backbone geometry.
\newblock In \emph{NeurIPS 2025}. arXiv:2503.05738v2, 2025.

\bibitem{yim2023framediff}
J.\ Yim et~al.
\newblock {SE(3)} diffusion model with application to protein backbone generation.
\newblock In \emph{ICML 2023}, 2023.

\bibitem{yim2023frameflow}
J.\ Yim et~al.
\newblock Fast protein backbone generation with {SE(3)} flow matching.
\newblock \emph{arXiv:2310.05297}, 2023.

\end{thebibliography}

\end{document}
